\section*{Task 3: Full-rank recovery}

\paragraph{Goal.}
Check that the main theorem recovers the global-minimum orbit picture when
\(r=d_L\).

\paragraph{Specialization.}
If \(r=d_L\), then
\[
P_Q=0,
\qquad
W_0=\Sigma_{XY}P_Q=0,
\qquad
G_d^{W_0}=G_d.
\]
The closing-arrow rank condition becomes
\[
\rho_{L,L+1}=0.
\]
Hence no wrapped interval can occur, so the admissible lifted rank patterns are
exactly the ordinary rank patterns of the equioriented chain
\[
0\to 1\to \cdots \to L.
\]

\paragraph{Resulting orbit classification.}
The main theorem therefore reduces to
\[
H_S\backslash \mathcal C_{d,S}^{d_L}
\xrightarrow{\sim}
\widetilde{\mathcal R}_{d,S}^{\,d_L},
\]
where the right-hand side is precisely the chain rank-pattern set classifying
the \(G_d\)-orbits in the full-rank determinantal locus.

\paragraph{Interpretation.}
This is the direct weight-space recovery of the global-minimum geometry studied
by Simon Pepin Lehalleur: the saddle-pattern theorem contains the global-minimum
fiber as its extremal full-rank case.

\paragraph{Status.}
Completed. The weight-space theorem is compatible with the known full-rank
orbit picture.
